% !TEX program = lualatex
% Generated by scripts/build_module_practice.py - do not edit by hand.
\DocumentMetadata{
  pdfstandard=UA-2,
  pdfversion=2.0,
  lang=en-US,
  tagging=on
}
\documentclass[11pt]{article}
\usepackage[letterpaper, margin=0.8in, top=0.6in, bottom=0.6in]{geometry}
\usepackage{fontspec}
\setmainfont{Fira Sans}
\usepackage{amsmath}
\usepackage{amssymb}
\usepackage{booktabs}
\usepackage{longtable}
\usepackage{array}
\usepackage{calc}
\usepackage{etoolbox}
\usepackage{xcolor}
\definecolor{accent}{HTML}{BF5700}
\usepackage{titlesec}
\titleformat{\section}{\large\bfseries\color{accent}}{}{0pt}{}
\titlespacing*{\section}{0pt}{14pt}{4pt}
\titleformat{\subsection}{\normalsize\bfseries}{}{0pt}{}
\titlespacing*{\subsection}{0pt}{10pt}{2pt}
\usepackage{hyperref}
\hypersetup{pdftitle={ECON 201 Practice Problems: The Economist's Toolkit}, pdfauthor={Div Bhagia}, hidelinks}
\pagestyle{plain}
\setlength{\parindent}{0pt}
\setlength{\parskip}{4pt}
\begin{document}
{\LARGE\bfseries\color{accent} Practice Problems: The Economist's Toolkit}\hfill{\large ECON 201}
\vspace{4pt}


\section{Lecture 3: Making economic decisions: opportunity cost, rents, and incentives}\label{lecture-3-making-economic-decisions-opportunity-cost-rents-and-incentives}

\subsection{1. Adopting a new technology}\label{adopting-a-new-technology}

Lena runs a small print shop. Her old press does the job, but it takes three operators. A supplier offers to lease her a digital press that needs one. Either way the shop sells the same work: \textbf{\$150,000 a year}.

\begin{itemize}
\tightlist
\item
  \textbf{Keep the old press.} Wages, ink, and maintenance come to \$120,000 a year.
\item
  \textbf{Lease the new press.} The lease is \$25,000 a year, and wages, ink, and maintenance fall to \$80,000. Total: \$105,000 a year.
\end{itemize}

\textbf{(a)} Make a table with a column for each option, and work out three rows: benefit, direct cost, and net benefit.

\textbf{(b)} What is the opportunity cost of each option?

\textbf{(c)} What is the economic cost of each option?

\textbf{(d)} Which option should Lena choose? What is her economic rent from that option?

\textbf{(e)} Two years later every print shop in town has the new press, and competition has pushed prices down: the same work now brings in \$120,000 a year with either press. Work out the net benefit of each option again. Is leasing still the right choice? What has happened to the extra profit Lena earned by moving first?

\subsection{2. Is the shop really profitable?}\label{is-the-shop-really-profitable}

Lena's cousin looks at her books from problem 1 and says she cannot lose: either press turns a profit. But suppose Lena could close the shop and earn \$40,000 a year managing another print shop.

\textbf{(a)} With the old press, the shop nets \$30,000 a year. Was running it really profitable?

\textbf{(b)} With the new press it nets \$45,000 a year. Is it profitable now, and by how much?

\textbf{(c)} In one or two sentences: why do the books and the economist disagree about the same shop?

\textbf{(d)} With the old press, the numbers say the job wins. Is there any justification for Lena keeping the shop anyway?

\subsection{3. Multiple choice}\label{multiple-choice}

\medskip

\textbf{1.} An airline sells a last-minute seat on an otherwise empty plane for \$25. The opportunity cost to the airline of filling that seat is:

a) The average cost per seat for the flight.

b) The full-fare ticket price.

c) Close to zero, since the plane flies either way and the seat would otherwise be empty.

d) Negative, since the passenger may buy a drink on board.

\medskip

\textbf{2.} You win a free ticket to see Band A tonight, and it cannot be resold. Band B plays at the same time; a Band B ticket costs \$40, and an evening with Band B is worth \$50 to you. Apart from the tickets, neither concert has any other cost. What is the opportunity cost of using your free ticket to see Band A?

a) \$0

b) \$10

c) \$40

d) \$50

\section{Lecture 4: Why do markets matter? Specialization and comparative advantage}\label{lecture-4-why-do-markets-matter-specialization-and-comparative-advantage}

\subsection{1. Pablo and Lin strike different deals}\label{pablo-and-lin-strike-different-deals}

From the lecture: if Pablo and Lin each spent their whole Sunday on just one food, they could produce:

\begin{longtable}[]{@{}lcc@{}}
\toprule\noalign{}
& Loaves of bread & Pounds of cheese \\
\midrule\noalign{}
\endhead
\bottomrule\noalign{}
\endlastfoot
Pablo & 12 & 6 \\
Lin & 2 & 4 \\
\end{longtable}

On their own, each spends 75\% of the day on bread, so Pablo produces 9 loaves and 1.5 pounds of cheese, and Lin produces 1.5 loaves and 1 pound. When they specialize, Pablo bakes 12 loaves and Lin makes 4 pounds of cheese.

\textbf{(a)} Suppose the price of cheese is 1.5 loaves per pound, and Lin sells Pablo 2 pounds. What does each consume, and is each better off than on their own?

\textbf{(b)} In the lecture, the deal was 2 pounds for 2 loaves, a price of 1 loaf per pound. Compare that with the 1.5-loaf price in part (a). Whose share of the gains from trade grew?

\textbf{(c)} Now suppose the price falls to a quarter of a loaf per pound. What happens?

\subsection{2. Two food trucks}\label{two-food-trucks}

Luis and Rosa each run a food truck. If each spent a whole workday on just one product, they could make:

\begin{longtable}[]{@{}lcc@{}}
\toprule\noalign{}
& Tacos & Burritos \\
\midrule\noalign{}
\endhead
\bottomrule\noalign{}
\endlastfoot
Luis & 40 & 8 \\
Rosa & 12 & 6 \\
\end{longtable}

Each can also split the day, with output in proportion to time.

\textbf{(a)} Find each person's opportunity cost of a burrito and of a taco.

\textbf{(b)} Who has the absolute advantage in each product? Who has the comparative advantage in each product?

\textbf{(c)} On their own, Luis spends three quarters of the day on tacos and Rosa spends two thirds of the day on tacos. How much does each produce, and what are the totals?

\textbf{(d)} Now each specializes completely in the product in which they have the comparative advantage. What is produced in total, and how does it compare with part (c)?

\textbf{(e)} They agree on a deal: Rosa sells Luis 3 burritos for 9 tacos. What does each of them end up with, and is each better off than in part (c)?

\textbf{(f)} Over what range of prices, in tacos per burrito, would both gain from trade?

\subsection{3. Ana and Ben at the campus bakery}\label{ana-and-ben-at-the-campus-bakery}

Ana and Ben work at a campus bakery. If each spent a whole hour on just one task, they could produce:

\begin{longtable}[]{@{}lcc@{}}
\toprule\noalign{}
& Frosted cupcakes & Pots of coffee \\
\midrule\noalign{}
\endhead
\bottomrule\noalign{}
\endlastfoot
Ana & 12 & 3 \\
Ben & 4 & 4 \\
\end{longtable}

\textbf{(a)} Who has the absolute advantage in each task?

\textbf{(b)} Who has the comparative advantage in each task?

\textbf{(c)} Suppose Ben sells coffee to Ana at a price measured in cupcakes per pot. Over what range of prices would both of them gain from the trade?

\subsection{4. Multiple choice}\label{multiple-choice}

\medskip

\textbf{1.} Maya paid \$45 for a concert ticket that is non-refundable and cannot be resold. On the day of the show she is exhausted and would much rather stay home. Which reasoning is economically sound?

a) She should go: otherwise the \$45 is wasted.

b) The \$45 is gone whether she goes or not; she should simply compare how much she would enjoy the concert against how much she values a night of rest.

c) She should go only if the ticket cost more than the value she places on a night at home.

d) Since the ticket is already paid for, going tonight costs her nothing, so she should go.

\medskip

\textbf{2.} In an hour, Priya can string 20 bracelets or 5 necklaces. Her opportunity cost of one necklace, in terms of bracelets, is:

a) A quarter of a bracelet.

b) 4 bracelets.

c) 5 bracelets.

d) 20 bracelets.

\medskip

\textbf{3.} Alex and Jose live on separate islands. If each spends all of their time on one good, Alex can produce 10,000 melons or 15,000 oranges, and Jose can produce 15,000 melons or 30,000 oranges. Select all the correct statements.

a) Jose has an absolute advantage in both goods.

b) Jose has an absolute advantage in oranges, and Alex has an absolute advantage in melons.

c) Jose's opportunity cost of producing an orange is higher than Alex's.

d) Alex has a comparative advantage in melons.

\medskip

\textbf{4.} Suppose Alex and Jose both specialize in the good in which they have the comparative advantage, and they trade at a price measured in melons per orange. Which statement is correct?

a) Alex specializes in oranges, and a higher price favors Alex.

b) Jose specializes in oranges, and a higher price favors Jose.

c) Jose specializes in oranges, and a higher price favors Alex.

d) The price does not affect how the gains from trade are shared.

\medskip

\textbf{5.} A country has an absolute advantage in every good it produces. It follows that:

a) It cannot gain from trading with anyone.

b) It has a comparative advantage in every good as well.

c) It still has a comparative advantage in some goods and not others, and can gain by trading.

d) Its trading partners would lose from trading with it.

\medskip

\textbf{6.} Ravi's opportunity cost of a kilo of coffee is 2 kilos of rice. Sofia's is 5 kilos of rice. At which price for a kilo of coffee could both of them gain from trading?

a) 1 kilo of rice.

b) 2 kilos of rice.

c) 3 kilos of rice.

d) 6 kilos of rice.

\medskip

\textbf{7.} From the lecture: suppose Taylor Swift types faster than any assistant she could hire. Should she answer her own fan mail?

a) Yes, she has an absolute advantage in typing.

b) Yes, paying an assistant is a direct cost, while typing it herself is free.

c) No, because the time she would spend typing is worth more than what an assistant would cost her.

d) It does not matter, since either way the mail gets answered.

\medskip

\textbf{8.} The United States could produce clothing using fewer worker-hours per shirt than many of the countries it imports clothing from. Why might importing still make sense?

a) The exporting countries must secretly have an absolute advantage.

b) American labor and capital have a high opportunity cost: hours spent sewing shirts would give up more valuable output in industries like aircraft and software.

c) Imported goods are always cheaper than domestically produced ones.

d) The United States lacks the technology to produce clothing.

\medskip

\textbf{9.} Which of the following are downsides of specialization? Select all that apply.

a) Total output falls when people specialize.

b) Specializing makes you dependent on others, so a disruption to trade can leave you badly exposed.

c) A specialized producer is vulnerable to a fall in demand or price for the one thing they make.

d) Specialization benefits only the more productive party.

\end{document}
